% ============================================================
% Section 175: 对称陀螺的能级及稍不对称时的修正
% ============================================================
\section{量子力学：对称陀螺的能级及稍不对称时的修正}
\label{sec:symmetric-top-asymmetry}

\begin{modelbox}[题目：对称陀螺的能级及稍不对称时的修正]
\begin{enumerate}
    \item 给出惯量主矩 $I_1=I_2=I\neq I_3$ 的对称陀螺的所有能级；
    \item 稍不对称陀螺满足 $I_1-I_2=\Delta\neq0$，$I_1+I_2=2I$，且 $\Delta/(2I)\ll1$。计算 $J=0$ 和 $J=1$ 能量到 $O(\Delta)$ 量级。
\end{enumerate}
\end{modelbox}

\begin{tipbox}[考点快速分析]
\begin{itemize}
    \item \textbf{对称陀螺}：$H_0=\dfrac{\mathbf{J}^2}{2I}+\dfrac12\bigl(\dfrac{1}{I_3}-\dfrac{1}{I}\bigr)J_z^2$，$J_z$ 为体轴角动量分量
    \item \textbf{本征态}：$|Jm\rangle$（$m$ 为体轴 $z$ 分量量子数），能级依赖于 $J$ 和 $m$
    \item \textbf{不对称微扰}：$H'=-\dfrac{\Delta}{8I^2}(J_+^2+J_-^2)$，选择定则 $\Delta m=\pm2$
    \item \textbf{简并子空间}：$J=1,\;m=\pm1$ 需对角化 $2\times2$ 有效矩阵
\end{itemize}
\end{tipbox}

\begin{derivationbox}[标准解题过程]
\textbf{1. 对称陀螺的能级}

固连坐标系中，$I_1=I_2=I$：
\[
H=\frac{1}{2I}(J_x^2+J_y^2)+\frac{1}{2I_3}J_z^2=\frac{\mathbf{J}^2}{2I}+\frac12\Bigl(\frac{1}{I_3}-\frac{1}{I}\Bigr)J_z^2.
\]

$H$ 与 $\mathbf{J}^2$、$J_z$ 对易，本征态 $|Jm\rangle$（$m$ 为体轴 $z$ 分量）。能级
\begin{equation}
\boxed{E_{Jm}=\frac{\hbar^2}{2I}J(J+1)+\frac{\hbar^2}{2}\Bigl(\frac{1}{I_3}-\frac{1}{I}\Bigr)m^2,}
\end{equation}
其中 $J=0,1,2,\ldots$，$m=-J,\ldots,J$。

\textbf{2. 稍不对称陀螺的修正}

设 $I_1=I+\Delta/2$，$I_2=I-\Delta/2$，$|\Delta|\ll I$。展开到一阶：
\[
H=\frac{J_x^2}{I+\Delta/2}+\frac{J_y^2}{I-\Delta/2}+\frac{J_z^2}{2I_3}\approx H_0-\frac{\Delta}{4I^2}(J_x^2-J_y^2)=H_0+H'.
\]

利用 $J_\pm=J_x\pm iJ_y$，有 $J_x^2-J_y^2=\frac12(J_+^2+J_-^2)$，故
\[
H'=-\frac{\Delta}{8I^2}(J_+^2+J_-^2).
\]

\textit{$J=0$}：唯一态 $|0,0\rangle$，$J_\pm^2|0,0\rangle=0$，故
\[
E_{00}'=0.
\]

\textit{$J=1,\;m=0$}：$J_\pm^2|1,0\rangle$ 产生 $|1,\pm2\rangle$（不存在），矩阵元为零，故
\[
E_{10}'=E_{10}=\frac{\hbar^2}{I}.
\]

\textit{$J=1,\;m=\pm1$}：子空间 $\{|1,1\rangle,|1,-1\rangle\}$ 中
\[
J_+^2|1,-1\rangle=2\hbar^2|1,1\rangle,\quad J_-^2|1,1\rangle=2\hbar^2|1,-1\rangle.
\]

$H'$ 矩阵为 $-\dfrac{\Delta\hbar^2}{4I^2}\begin{pmatrix}0&1\\1&0\end{pmatrix}$。对角化得本征值 $\pm\dfrac{\Delta\hbar^2}{4I^2}$。因此
\begin{equation}
\boxed{E_{1,\pm1}^{(\pm)}=\frac{\hbar^2}{2I}+\frac{\hbar^2}{2I_3}\pm\frac{\Delta\hbar^2}{4I^2}.}
\end{equation}
\end{derivationbox}

\begin{formulabox}[答案汇总]
\begin{center}
\begin{tabular}{ll}
\toprule
\textbf{结论} & \textbf{结果} \\
\midrule
对称陀螺能级 & $E_{Jm}=\dfrac{\hbar^2}{2I}J(J+1)+\dfrac{\hbar^2}{2}\bigl(\dfrac{1}{I_3}-\dfrac{1}{I}\bigr)m^2$ \\
$J=0$ 修正 & $E_{00}'=0$ \\
$J=1,m=0$ 修正 & $E_{10}'=E_{10}=\hbar^2/I$ \\
$J=1,m=\pm1$ 分裂 & $E_{1,\pm1}^{(\pm)}=\dfrac{\hbar^2}{2I}+\dfrac{\hbar^2}{2I_3}\pm\dfrac{\Delta\hbar^2}{4I^2}$ \\
\bottomrule
\end{tabular}
\end{center}
\end{formulabox}

\begin{insightbox}[物理图像]
\begin{itemize}
    \item 对称陀螺的能级公式是分子光谱学的基础。对于 $I_3<I$（扁球陀螺，如盘状分子），$m\neq0$ 的态能量升高；对于 $I_3>I$（长球陀螺，如杆状分子），$m\neq0$ 的态能量降低。
    \item 稍不对称性（$\Delta\neq0$）是理解真实分子转动能级精细结构的起点。微波光谱中可观测到由不对称性导致的能级分裂，进而反推分子的惯量张量。
    \item 本题中 $J=1,m=\pm1$ 的简并解除对应于绕不同体轴转动的能量差异：$|1,1\rangle\pm|1,-1\rangle$ 组合分别对应绕 $x$ 轴和 $y$ 轴的转动，因 $I_1\neq I_2$ 而能量不同。
\end{itemize}
\end{insightbox}
